% =====================================================================
%  PAGES LIMINAIRES : titre, promesse, mode d'emploi, feuille de route
% =====================================================================
\begin{titlepage}
\thispagestyle{empty}
\begin{center}

\vspace*{0.6cm}
{\color{orange}\rule{\linewidth}{2.5pt}}\\[0.9cm]

{\sffamily\bfseries\color{bleuNuit}\fontsize{34}{40}\selectfont ANGLAIS\par}
\vspace{0.35cm}
{\sffamily\bfseries\color{bleuNuit}\fontsize{21}{26}\selectfont Le grand livre d'auto-formation\par}
\vspace{0.5cm}
{\color{gris}\large\itshape 12 mois pour parler anglais\\ comme vous parlez français\par}

\vspace{0.7cm}
{\color{orange}\rule{0.45\linewidth}{1.5pt}}

\vspace{0.9cm}
{\sffamily\large\color{bleuNuit}\bfseries Méthode complète pour francophones\par}
\vspace{0.25cm}
{\large Niveau de départ A1/A2 \quad\ra\quad Objectif B2 (conversation courante)\par}

\vspace{1.1cm}

\begin{minipage}{0.92\linewidth}
\centering
\color{bleuNuit}\large\bfseries
Ce que contient ce livre
\end{minipage}

\vspace{0.35cm}

\begin{minipage}{0.94\linewidth}
\small
\begin{tabular}{@{}p{7.3cm}p{7.3cm}@{}}
\ok Un plan de 12 mois, semaine par semaine & \ok Plus de 3\,000 mots et 800 phrases prêtes à l'emploi \\[4pt]
\ok Toute la grammaire, expliquée en français & \ok 40 situations de conversation, dialogues modèles et jeux de rôle \\[4pt]
\ok La prononciation : 44 sons, accent, intonation & \ok 150 \en{phrasal verbs}, 120 expressions idiomatiques \\[4pt]
\ok Les 100 erreurs typiques des francophones & \ok Des centaines d'exercices, tous corrigés \\[4pt]
\ok La méthode d'écoute, de lecture et d'écriture & \ok Tests de niveau et grilles d'auto-évaluation \\
\end{tabular}
\end{minipage}

\vfill

{\color{gris}\small Édition 1.0 --- Document LaTeX librement modifiable et imprimable.\par}
\vspace{0.3cm}
{\color{orange}\rule{\linewidth}{2.5pt}}

\end{center}
\end{titlepage}

% ---------------------------------------------------------------------
\chapter*{Avant de commencer}
\addcontentsline{toc}{chapter}{Avant de commencer}
\markboth{Avant de commencer}{Avant de commencer}

Vous avez déjà des notions d'anglais. Vous comprenez peut-être un peu, vous lisez des mots
isolés, mais dès qu'il faut \emph{parler}, tout se bloque : les mots ne viennent pas, la
prononciation vous gêne, vous traduisez du français dans votre tête et vous arrivez deux
secondes trop tard.

Ce livre a été écrit pour exactement ce problème. Il ne s'agit pas d'un manuel scolaire de
plus, mais d'un \textbf{programme d'entraînement complet} : une méthode, un calendrier, des
contenus et des exercices, organisés pour qu'en douze mois vous passiez de « je connais
un peu l'anglais » à « je parle anglais, je comprends ce qu'on me dit et je réponds sans
hésiter ».

\begin{memo}[La promesse de ce livre]
Si vous suivez le plan du chapitre 2 --- \textbf{90 minutes par jour, six jours sur sept} ---
pendant douze mois, vous aurez accumulé environ \textbf{470 heures d'entraînement dirigé}
et \textbf{3\,000 heures de contact} avec l'anglais si vous ajoutez l'immersion passive
recommandée. C'est le volume qui sépare un francophone débutant d'un locuteur \textbf{B2
solide}, capable de tenir une conversation de 30 minutes avec un anglophone, de suivre un
film sans sous-titres français et de travailler en anglais.
\end{memo}

\section*{Pourquoi un francophone peut réussir vite}
Trois bonnes nouvelles, à lire les jours de découragement :

\begin{enumerate}
  \item \textbf{L'anglais et le français partagent environ 30 à 45\,\% de leur vocabulaire.}
  Depuis 1066, le français a massivement irrigué l'anglais : \en{government}, \en{language},
  \en{important}, \en{possible}, \en{different}, \en{information}, \en{nature}, \en{culture},
  \en{justice}, \en{village}, \en{table}, \en{fruit}, \en{hotel}. Vous connaissez déjà des
  milliers de mots anglais sans le savoir. Le chapitre 12 vous apprend à exploiter ce trésor.
  \item \textbf{La grammaire anglaise est plus simple que la grammaire française.} Pas de
  genre sur les noms, pas d'accord des adjectifs, pas de subjonctif à mémoriser, une seule
  forme de verbe au présent (\en{-s} à la 3\textsuperscript{e} personne), pas de vouvoiement.
  Ce qui est difficile en anglais, ce sont \emph{les temps} (le \en{present perfect} !)
  et \emph{la prononciation}. C'est précisément là que ce livre insiste.
  \item \textbf{L'anglais est partout.} Musique, séries, football, technologie, réseaux
  sociaux, religion, commerce : vous pouvez trouver de l'anglais gratuit toute la journée.
  L'immersion ne vous coûtera presque rien --- seulement de la discipline.
\end{enumerate}

\section*{Les trois pièges qui font échouer (et comment les éviter)}
\begin{tableau}[Ce qui fait échouer 90 \% des autodidactes]
\begin{tabularx}{\linewidth}{@{}p{3.9cm}X@{}}
\entete{Erreur} & \entete{Correction apportée par ce livre} \tabularnewline
\hline
\textbf{Étudier la grammaire sans jamais parler} &
Chaque unité se termine par une \textbf{mission orale} : dire 10 phrases à voix haute,
se raconter quelque chose, imiter un dialogue. La parole est entraînée dès le jour 1. \tabularnewline
\hline
\textbf{Apprendre des listes de mots isolés} &
Vous apprendrez des \textbf{blocs} (\en{I'm looking forward to it}) et non des mots
seuls, avec la prononciation, et vous les reverrez grâce à la répétition espacée
(chapitre 3). \tabularnewline
\hline
\textbf{Recommencer dix fois le niveau débutant} &
Un \textbf{calendrier daté} avec des jalons : mois 1, mois 3, mois 6, mois 12. Vous savez
toujours où vous en êtes, et les tests du chapitre 19 le prouvent. \tabularnewline
\end{tabularx}
\end{tableau}

\section*{Comment lire ce livre}
Ce livre utilise un code couleur constant. En voici la légende --- gardez cette page en tête
pendant toute votre année :

\begin{regle}
Une règle de grammaire ou de prononciation, formulée simplement, avec sa structure.
\end{regle}

\begin{piege}
Un endroit où les francophones se trompent presque toujours. Ce sont les encadrés les plus
rentables du livre : chaque piège évité, c'est une faute que vous ne ferez plus jamais.
\end{piege}

\begin{astuce}
Un raccourci de mémorisation, une astuce de prononciation, un moyen de gagner du temps.
\end{astuce}

\begin{mission}
Ce que vous devez \textbf{faire} aujourd'hui, à voix haute. Un livre de langue ne s'apprend
pas avec les yeux : une mission non faite est une leçon non apprise.
\end{mission}

\begin{memo}
La synthèse de l'unité : à relire la veille d'un test ou dans le bus.
\end{memo}

\begin{exo}[X --- modèle]
Les exercices sont numérotés par unité. Le \textbf{corrigé suit immédiatement} l'exercice
(encadré vert), pour que vous puissiez travailler seul sans attendre. Cachez-le avec une
feuille si vous préférez chercher d'abord.
\end{exo}

\begin{corrige}
Voici à quoi ressemble un corrigé.
\end{corrige}

% Conventions d'écriture
\begin{tableau}[Conventions typographiques]
\begin{tabularx}{\linewidth}{@{}p{5.4cm}X@{}}
\entete{Ce que vous voyez} & \entete{Ce que cela veut dire} \tabularnewline
\hline
\en{I have been working.} & Texte en \textbf{anglais} (à lire à voix haute). \tabularnewline
\hline
\textit{\color{gris}J'ai travaillé / je travaille depuis\ldots} & Traduction française, en petit et en gris. \tabularnewline
\hline
\pron{aïe have bin oueur-king} & \textbf{Prononciation approchée à la française}.
Indispensable les premiers mois ; passez ensuite à l'écoute réelle (chapitre 16). \tabularnewline
\hline
\trou \quad ou \quad \trou[1.2cm] & Trou à compléter dans un exercice. \tabularnewline
\hline
\en{*I have 25 years.} & L'astérisque signale une phrase \textbf{incorrecte}, à ne pas imiter. \tabularnewline
\hline
\en{BrE} / \en{AmE} & Britannique / américain. Les deux sont donnés quand ils diffèrent. \tabularnewline
\end{tabularx}
\end{tableau}

\begin{astuce}[Le mot \en{to be} ou \en{to have} ?]
Dans ce livre, tout exemple anglais est suivi de sa traduction. Ne lisez \textbf{jamais}
un exemple sans le dire à voix haute au moins une fois. Votre bouche doit s'habituer, pas
seulement vos yeux.
\end{astuce}

% ---------------------------------------------------------------------
\chapter*{Feuille de route en un coup d'\oe il}
\addcontentsline{toc}{chapter}{Feuille de route en un coup d'\oe il}
\markboth{Feuille de route}{Feuille de route}

Voici votre année, résumée sur une page. Le détail jour par jour se trouve au chapitre 2.

\begin{tableau}[Les 12 mois --- objectifs, contenu, jalon de validation]
\footnotesize
\begin{longtable}{@{}p{1.7cm}p{2.1cm}p{6.3cm}p{4.6cm}@{}}
\entete{Mois} & \entete{Niveau visé} & \entete{Ce que vous travaillez} & \entete{Vous devez savoir faire} \tabularnewline
\endfirsthead
\entete{Mois} & \entete{Niveau visé} & \entete{Ce que vous travaillez} & \entete{Vous devez savoir faire} \tabularnewline
\endhead
\hline
\textbf{1} & A1 --- démarrage &
Sons de l'anglais, accent tonique, \en{be}/\en{have}, présent simple, se présenter,
300 mots, écoute quotidienne 15 min &
Se présenter en 2 min, dire 200 phrases de base, épeler son nom \tabularnewline
\hline
\textbf{2} & A1 --- consolidation &
Questions, négation, articles, pluriels, \en{there is}, prépositions de lieu/temps,
vie quotidienne, 600 mots &
Raconter sa journée, poser 30 questions types \tabularnewline
\hline
\textbf{3} & A1\ra A2 &
Passé simple (verbes irréguliers), \en{past continuous}, futur \en{going to}/\en{will},
comparatifs, 900 mots, 1\textsuperscript{er} dialogue réel &
Raconter un souvenir, décrire un projet, tenir 5 min de conversation guidée \tabularnewline
\hline
\textbf{4} & A2 &
\en{Present perfect} (le grand chapitre), \en{for}/\en{since}, \en{used to},
modaux de base, ville/transports/santé, 1\,200 mots &
Parler de son expérience, réserver, demander son chemin, chez le médecin \tabularnewline
\hline
\textbf{5} & A2 &
Tous les temps en revue, voix passive simple, \en{should}/\en{must}/\en{have to},
travail et métiers, 1\,500 mots, lecture d'un premier roman simplifié &
Décrire un emploi, une règle, un processus ; lire un texte A2 sans dictionnaire \tabularnewline
\hline
\textbf{6} & A2\ra B1 &
\textbf{Bilan à mi-parcours.} Conditionnels 0/1/2, relatifs, gérondif/infinitif,
\en{phrasal verbs} (1\textsuperscript{re} série), 1\,800 mots &
Donner son opinion, nuancer, comprendre une série avec sous-titres anglais \tabularnewline
\hline
\textbf{7} & B1 &
Conditionnel 3, discours rapporté, \en{wish}, inversion, quantifieurs,
médias et actualités, 2\,100 mots &
Raconter une anecdote complexe, rapporter des propos, débattre 10 min \tabularnewline
\hline
\textbf{8} & B1 &
Voix passive complète, \en{have something done}, \en{cleft sentences},
modaux du passé, voyages/aéroport/hôtel, 2\,400 mots &
Se débrouiller en voyage de A à Z, présenter un exposé de 5 min \tabularnewline
\hline
\textbf{9} & B1\ra B2 &
Connecteurs avancés, collocations, expressions idiomatiques, registre
(formel/familier), 2\,700 mots, écriture d'e-mails et de lettres &
Écrire un e-mail professionnel, adapter son registre, comprendre un podcast clair \tabularnewline
\hline
\textbf{10} & B2 &
Anglais professionnel : réunion, entretien d'embauche, téléphone, négociation,
présentation ; 3\,000 mots &
Passer un entretien, animer une réunion, téléphoner sans panique \tabularnewline
\hline
\textbf{11} & B2 &
Fluidité : débit, \emph{fillers}, reformulation, humour, accents (britannique,
américain, africain, indien), 30 min de conversation libre par semaine &
Parler 20 min sans préparation sur un sujet inconnu \tabularnewline
\hline
\textbf{12} & B2\ra C1 &
Révision générale, test blanc complet, certification éventuelle (TOEIC/IELTS),
consolidation des points faibles &
Tenir une conversation d'une heure, regarder un film sans sous-titres \tabularnewline
\end{longtable}
\end{tableau}

\begin{mission}[Avant même de commencer]
\begin{enumerate}
  \item Passez le \textbf{test de positionnement} du chapitre 19 (40 min). Il vous dira si
  vous démarrez au mois 1 ou au mois 3.
  \item Choisissez votre \textbf{créneau quotidien} et notez-le : \trou[3cm] h \ra \trou[2cm] h.
  \item Installez \textbf{une application de répétition espacée} (Anki, gratuite) et
  \textbf{un dictionnaire} (Cambridge ou WordReference).
  \item Mettez votre \textbf{téléphone en anglais} aujourd'hui (Réglages \ra Langue).
  Cinq minutes de courage, un an de bénéfices.
  \item Dites à voix haute, maintenant : \en{``In one year, I will speak English every day
  without fear.''} --- \textit{Dans un an, je parlerai anglais tous les jours sans peur.}
\end{enumerate}
\end{mission}
